% !TEX root = ../projekt.tex
% =====================================================================
%  PRÍLOHA A — OBSAH PAMÄŤOVÉHO MÉDIA
% =====================================================================
\chapter{Obsah priloženého pamäťového média}
\label{app:medium}

\rozsah{0,5--1 strana. Povinná príloha — over si aktuálne požiadavky
fakulty na štruktúru odovzdávaného archívu.}

\pokyn{Vypíš stromovú štruktúru archívu a pri každom adresári jednu vetu,
čo obsahuje. Odporúčaná štruktúra:}

\begin{verbatim}
/
|-- text/        zdrojove subory prace (LaTeX) a vysledne PDF
|-- src/         evaluacny nastroj (zdrojove kody, requirements)
|-- rtl/         VHDL zdrojove kody rozsirenej RSS komponenty
|-- data/        vzorka vstupnych dat (alebo odkaz, ak sa nesmu zverejnit)
|-- results/     vysledne CSV subory a vygenerovane grafy
|-- scripts/     skripty na zopakovanie experimentov
`-- README.md    navod na spustenie
\end{verbatim}

\pozor{Pred odovzdaním uprac repozitár:
\begin{itemize}
  \item \texttt{README.md} je zastaraný (37\,B tuple, 64\,B kľúč, staré
        stĺpce a staré algoritmy) — musí zodpovedať finálnemu stavu,
  \item zmaž zvyšky \texttt{results/t.csv} a \texttt{resultsoverload\_comparison.png},
  \item \texttt{run.sh} odkazuje na neexistujúci \texttt{data/odid1.csv},
  \item \texttt{analyze.py:191} má pokazený nadpis grafu,
  \item oprav preklep \texttt{thresshold\_sum} $\rightarrow$ \texttt{threshold\_sum}.
\end{itemize}
Oponent si archív otvorí a skúsi spustiť \texttt{run.sh}. Ak to spadne
na prvom riadku, poznačí si to.}

\pokyn{Ak sa dáta CESNET nesmú zverejniť, napíš to sem explicitne a priloži
aspoň malú anonymizovanú vzorku alebo generátor syntetických dát, aby sa
nástroj dal spustiť aj bez nich.}
